\documentclass[12pt,a4paper]{report}
\usepackage[utf8]{inputenc}
\usepackage{graphicx}
\usepackage{geometry}
\usepackage{setspace}
\usepackage{array}
\usepackage{multirow}

% Page margins matching standard report requirements
\geometry{top=1in, bottom=1in, left=1.25in, right=1in}

\begin{document}

% --- 1. TITLE PAGE ---
\begin{titlepage}
    \begin{center}
        \setstretch{1.5}
        \textbf{SAREE BLOUSE MATCHER: \\ COMPUTER VISION BASED MATCHING SYSTEM} \\
        \vspace{0.8cm}
        \textbf{MINI PROJECT REPORT} \\
        \vspace{1.2cm}
        SUBMITTED IN PARTIAL FULFILLMENT OF THE REQUIREMENTS \\ FOR THE AWARD OF THE DEGREE OF \\
        \vspace{1.2cm}
        \textbf{MASTER OF ENGINEERING} \\
        IN \\
        \textbf{COMPUTER SCIENCE AND ENGINEERING} \\
        \vfill
        SUBMITTED BY \\
        \textbf{NITHISH KUMAR R \hspace{1cm} (REG.NO. 2536020005)} \\
        \vfill
        \textbf{ANNAMALAI} \includegraphics[height=1.5cm]{logo.png} \textbf{UNIVERSITY} \\
        \vspace{0.3cm}
        \textbf{FACULTY OF ENGINEERING \& TECHNOLOGY} \\
        \textbf{DEPARTMENT OF COMPUTER SCIENCE AND ENGINEERING} \\
        \textbf{ANNAMALAI UNIVERSITY ANNAMALAI NAGAR-608002} \\
        \vspace{0.5cm}
        \textbf{MAY – 2026}
    \end{center}
\end{titlepage}

% --- 2. BONAFIDE CERTIFICATE ---
\newpage
\begin{center}
    \textbf{ANNAMALAI} \includegraphics[height=1.2cm]{logo.png} \textbf{UNIVERSITY} \\
    \textbf{DEPARTMENT OF COMPUTER SCIENCE AND ENGINEERING} \\
    \vspace{1.5cm}
    \textbf{BONAFIDE CERTIFICATE}
\end{center}
\vspace{1.5cm}
\begin{spacing}{1.5}
This is to certify that this report titled \textbf{SAREE BLOUSE MATCHER:AI-POWERED OBJECT DETECTION SYSTEM} is the bonafide record of the work done by \textbf{NITHISH KUMAR R} (Reg.No:2536020005) in partial fulfilment of the requirement for the Degree of Master of Engineering (Computer Science and Engineering) during the year 2025-2027.
\end{spacing}

\vspace{2cm}
\noindent \textbf{PROJECT INCHARGE} \hfill \textbf{Dr. R. BHAVANI} \\
1. \null \hfill Professor and Head \\
\null \hfill Dept. of Computer Science and \\
\null \hfill Engineering \\
\null \hfill Annamalai University \\

\vspace{1.5cm}
\noindent \textbf{INTERNAL EXAMINER} \hfill \textbf{EXTERNAL EXAMINER} \\

\vfill
\noindent PLACE: ANNAMALAI NAGAR \\
DATE:

% --- 3. ACKNOWLEDGEMENT ---
\newpage
\begin{center}
    \textbf{ACKNOWLEDGEMENT}
\end{center}
\vspace{1cm}
\begin{spacing}{1.5}
I owe our Special Thanks to \textbf{Dr. R. BHAVANI, Professor and Head of the Department of Computer Science and Engineering}, for her encouraging and valuable counsel throughout this Project. \\

I would like to express my deep sense of gratitude and sincere thanks to \textbf{Dr. N. J. NALINI Professor, Department of Computer Science and Engineering, Project coordinator}, for her canalized guidance, dedication, constant inspiration, and encouragement. \\

Last, but not least, I would like to thank my family and teammates for their support and encouragement for the successful completion of the thesis work.
\end{spacing}

\vspace{2cm}
\begin{center}
    \textbf{NITHISH KUMAR R \quad (Reg. No: 2536020005)}
\end{center}
% --- 4. TABLE OF CONTENTS ---
\newpage
\begin{center}
    \textbf{TABLE OF CONTENTS}
\end{center}
\vspace{0.5cm}

\renewcommand{\arraystretch}{1.5} 

\noindent
\begin{tabular}{|>{\centering\arraybackslash}p{2.5cm}|>{\centering\arraybackslash}p{8.5cm}|>{\centering\arraybackslash}p{2.5cm}|}
\hline
\textbf{CHAPTER NO.} & \textbf{TITLE} & \textbf{PAGE NO.} \\ \hline
& \textbf{ABSTRACT} & VI \\ \hline
\textbf{1} & \textbf{INTRODUCTION} & 1 \\ \hline
& 1.1 Overview & 1 \\ \hline
& 1.2 Project Objectives & 2 \\ \hline
& 1.3 Problem Statement & 3 \\ \hline
& 1.4 Problems with Traditional Methods & 4 \\ \hline
& 1.5 Background and Motivation & 6 \\ \hline
& 1.6 Deep Learning and Computer Vision Approach & 8 \\ \hline
& 1.7 System Workflow & 10 \\ \hline
\textbf{2} & \textbf{SURVEY OF TECHNOLOGY} & 12 \\ \hline
& 2.1 Literature Survey & 12 \\ \hline
& 2.2 Technology Stack & 14 \\ \hline
\textbf{3} & \textbf{SYSTEM ANALYSIS AND REQUIREMENTS} & 18 \\ \hline
& 3.1 Problem Analysis & 18 \\ \hline
& 3.2 Proposed System & 20 \\ \hline
& 3.3 Software and Hardware Requirements & 24 \\ \hline
\textbf{4} & \textbf{SYSTEM DESIGN AND IMPLEMENTATION} & 26 \\ \hline
& 4.1 System Architecture & 26 \\ \hline
& 4.2 Implementation Modules & 28 \\ \hline
& 4.3 API Endpoints & 32 \\ \hline
\textbf{5} & \textbf{RESULTS AND DISCUSSION} & 35 \\ \hline
\textbf{6} & \textbf{CONCLUSIONS} & 40 \\ \hline
\textbf{7} & \textbf{REFERENCES} & 42 \\ \hline
\end{tabular}
% --- 5. ABSTRACT ---
\newpage
\begin{center}
    \textbf{ABSTRACT}
\end{center}
\vspace{0.5cm}
\begin{spacing}{1.2}
Fashion coordination, particularly in traditional Indian attire, poses significant challenges for individuals seeking to match blouses with sarees. This project presents an AI-powered object detection system called ”Saree Blouse Matcher” that automates the process of identifying and recommending matching blouses for sarees using computer vision and deep learning technologies. \\

The system employs OpenAI’s CLIP (Contrastive Language-Image Pre-training) model for visual feature extraction, generating 512-dimensional embeddings that capture both low-level visual attributes and high-level semantic concepts. A hybrid matching algorithm combines visual similarity (40\%), metadata-based rules (40\%), and learned patterns from user feedback (20\%) to achieve accurate recommendations. Color analysis is performed using K-means clustering with HSV-based background filtering to extract dominant color palettes from garment images. \\

The system architecture consists of a FastAPI backend server with RESTful endpoints, a responsive web-based frontend with HTML5/CSS3/JavaScript, and a continuous learning module (MatchTrainer) that stores user feedback to improve accuracy over time. The grid-based detection system enables identification of matching blouses in composite images containing multiple items, while real-time webcam support provides instant feedback for live detection scenarios. \\

Performance evaluation demonstrates that the system achieves 80.75\% accuracy initially, improving to 91.5\% after 20-30 user feedback iterations. The CLIP-based color detection method outperforms traditional RGB K-means clustering with 92\% accuracy compared to 68\%. Average processing time is 1.5 seconds per image on CPU, reducing to under 0.5 seconds with GPU acceleration. \\

\noindent \textbf{Keywords:} Computer Vision, Deep Learning, CLIP Model, Fashion Technology, Object Detection, Color Analysis, Hybrid Matching Algorithm, Traditional Attire, Real-time Detection, Continuous Learning.
\end{spacing}

\end{document}